\chapter{2012年浙江卷（理）}

\section{单选题}

\begin{problem}
设集合 \(A = \{x \mid 1 < x < 4\}\)，集合 \(B = \{x \mid x^2 - 2x - 3 \leqslant 0\}\)，则 \(A \cap (\complement_{\R} B) =\)（\quad）

\choices
    {\((1, 4)\)}
    {\((3, 4)\)}
    {\((1, 3)\)}
    {\((1,2)\) \(\cup\) \((3,4)\)}
\end{problem}

\begin{answer}
B
\end{answer}

\begin{solution}
由
\[
x^2-2x-3=(x-3)(x+1)
\]
得 \(B=[-1,3]\)，从而
\(\complement_{\R}B=(-\infty,-1)\cup(3,+\infty)\).
因此
\[
A\cap(\complement_{\R}B)=(1,4)\cap\bigl((-\infty,-1)\cup(3,+\infty)\bigr)=(3,4)
\]
故选第二项.
\end{solution}

\begin{problem}
已知 \(\i\) 是虚数单位，则 \( \frac{3+\i}{1-\i}= \)（\quad）

\choices
    {\(1-2\i\)}
    {\(2-\i\)}
    {\(2 + \i\)}
    {\(1 + 2\i\)}
\end{problem}

\begin{answer}
D
\end{answer}

\begin{solution}
分母有理化，得
\[
\frac{3+\i}{1-\i}=\frac{(3+\i)(1+\i)}{(1-\i)(1+\i)}
=\frac{2+4\i}{2}=1+2\i
\]
故选第四项.
\end{solution}

\begin{problem}
设 \(a \in \R\)，则“\(a = 1\)”是直线 \(l_1: ax + 2y - 1 = 0\) 与直线 \(l_2: x + (a + 1)y + 4 = 0\) 平行的（\quad）

\choices
    {充分不必要条件}
    {必要不充分条件}
    {充分必要条件}
    {既不充分也不必要条件}
\end{problem}

\begin{answer}
A
\end{answer}

\begin{solution}
两直线平行的充要条件为其法向量平行，即
\[
\begin{vmatrix}a&2\\1&a+1\end{vmatrix}=0
\]
所以
\[
a(a+1)-2=0\Longleftrightarrow(a-1)(a+2)=0
\]
即 \(a=1\) 或 \(a=-2\).
当 \(a=1\) 时，两直线分别为 \(x+2y-1=0\) 和 \(x+2y+4=0\)，确实平行；但平行还可以由 \(a=-2\) 得到. 因此“\(a=1\)”是充分不必要条件，故选第一项.
\end{solution}

\begin{problem}
把函数 \( y = \cos 2x + 1 \) 的图象上所有点的横坐标伸长到原来的 \(2\) 倍（纵坐标不变），然后向左平移 \(1\) 个单位长度，再向下平移 \(1\) 个单位长度，得到的图象是 （\quad）

\choices
    {\choicebitmap[width=0.15\paperwidth]{img/2012/zhejiang_science/q4_optA.png}}
    {\choicebitmap[width=0.15\paperwidth]{img/2012/zhejiang_science/q4_optB.png}}
    {\choicebitmap[width=0.15\paperwidth]{img/2012/zhejiang_science/q4_optC.png}}
    {\choicebitmap[width=0.15\paperwidth]{img/2012/zhejiang_science/q4_optD.png}}
\end{problem}

\begin{answer}
A
\end{answer}

\begin{solution}
横坐标伸长到原来的 \(2\) 倍后，函数变为 \(y=\cos x+1\). 向左平移 \(1\) 个单位长度后，函数变为 \(y=\cos(x+1)+1\)；再向下平移 \(1\) 个单位长度，得到
\[
y=\cos(x+1)
\]
该函数的周期为 \(2\pi\)，且在 \(x=-1\) 处取得最大值，在 \(x=\frac{\pi}{2}-1\) 处与 \(x\) 轴相交，所以图象为第一种.
\end{solution}

\begin{problem}
设 \(\bs{a}\)，\(\bs{b}\) 是两个非零向量（\quad）

\choices
    {若 \(|\bs{a} + \bs{b}| = |\bs{a}| - |\bs{b}|\)，则 \(\bs{a} \perp \bs{b}\)}
    {若 \(\bs{a} \perp \bs{b}\)，则 \(|\bs{a} + \bs{b}| = |\bs{a}| - |\bs{b}|\)}
    {若 \(|\bs{a} + \bs{b}| = |\bs{a}| - |\bs{b}|\)，则存在实数 \(\lambda\)，使得 \(\bs{b} = \lambda \bs{a}\)}
    {若存在实数 \(\lambda\)，使得 \(\bs{b} = \lambda \bs{a}\)，则 \(|\bs{a} + \bs{b}| = |\bs{a}| - |\bs{b}|\)}
\end{problem}

\begin{answer}
C
\end{answer}

\begin{solution}
若 \(|\bs{a}+\bs{b}|=|\bs{a}|-|\bs{b}|\)，则右端非负，且两边平方得
\[
|\bs{a}|^2+|\bs{b}|^2+2\bs{a}\cdot\bs{b}
=|\bs{a}|^2+|\bs{b}|^2-2|\bs{a}||\bs{b}|
\]
因此 \(\bs{a}\cdot\bs{b}=-|\bs{a}||\bs{b}|\)，这说明两向量反向共线，故存在实数 \(\lambda\)，使得 \(\bs{b}=\lambda\bs{a}\)，第三项正确. 例如取 \(\bs{b}=-\frac12\bs{a}\)，等式成立但两向量不垂直，故第一项错误；若两向量垂直，则 \(|\bs{a}+\bs{b}|=\sqrt{|\bs{a}|^2+|\bs{b}|^2}\)，一般不等于 \(|\bs{a}|-|\bs{b}|\)，故第二项错误；取 \(\lambda=1\) 时等式变为 \(2|\bs{a}|=0\)，故第四项错误.
严格地说，若两向量垂直且题设等式成立，则平方后有 \(|\bs{a}|^2+|\bs{b}|^2=(|\bs{a}|-|\bs{b}|)^2\)，从而 \(2|\bs{a}||\bs{b}|=0\)，与两向量非零矛盾.
\end{solution}

\begin{problem}
若从 \(1\)，\(2\)，\(3\)，\(\cdots\)，\(9\) 这 9 个整数中同时取 4 个不同的数，其和为偶数，则不同的取法共有（\quad）

\choices
    {\(60\text{ 种}\)}
    {\(63\text{ 种}\)}
    {\(65\text{ 种}\)}
    {\(66\text{ 种}\)}
\end{problem}

\begin{answer}
D
\end{answer}

\begin{solution}
这 9 个数中有 5 个奇数和 4 个偶数. 四个数的和为偶数，当且仅当取出的偶数个数为 \(0\)、\(2\) 或 \(4\). 因而取法数为
\[
\mathrm C_5^4\mathrm C_4^0+\mathrm C_5^2\mathrm C_4^2+\mathrm C_5^0\mathrm C_4^4
=5+60+1=66
\]
故选第四项.
\end{solution}

\begin{problem}
设 \(S_{n}\) 是公差为 \(d\)（\(d \neq 0\)） 的无穷等差数列 \(\{a_{n}\}\) 的前 \(n\) 项和，则下列命题错误的是（\quad）

\choices
    {若 \(d < 0\)，则数列 \(\{S_{n}\}\) 有最大项}
    {若数列 \(\{S_{n}\}\) 有最大项，则 \(d < 0\)}
    {若数列 \(\{S_{n}\}\) 是递增数列，则对任意 \(n \in \N^{*}\)，均有 \(S_{n} > 0\)}
    {若对任意 \(n \in \N^{*}\)，均有 \(S_n > 0\)，则数列 \(\{S_n\}\) 是递增数列}
\end{problem}

\begin{answer}
C
\end{answer}

\begin{solution}
由
\[
S_n-S_{n-1}=a_n=a_1+(n-1)d
\]
可知 \(\{S_n\}\) 的相邻项差构成等差数列.

当 \(d<0\) 时，\(S_n\) 是关于 \(n\) 的开口向下二次函数在正整数上的取值，且 \(S_n\to-\infty\)，所以有最大项，第一项正确. 若 \(d>0\)，则 \(S_n\to+\infty\)，不可能有最大项，因此有最大项时必有 \(d<0\)，第二项正确.

若 \(\{S_n\}\) 递增，则 \(S_n-S_{n-1}=a_n\geqslant0\) 对 \(n\geqslant2\) 成立，并不能推出 \(S_1=a_1>0\). 例如取 \(a_1=-1\)，\(d=2\)，则 \(S_1=-1\)，\(S_2=0\)，\(S_3=3\)，\(\ldots\)，数列 \(\{S_n\}\) 递增但 \(S_1<0\)，所以第三项错误.

若对任意 \(n\in\N^{*}\) 都有 \(S_n>0\)，则 \(S_1=a_1>0\). 此时不能有 \(d<0\)，否则 \(S_n\to-\infty\)；由于 \(d\ne0\)，故 \(d>0\)，从而所有 \(a_n>0\)，即 \(S_n\) 递增，第四项正确. 故选第三项.
\end{solution}

\begin{problem}
如图，\(F_{1}\)，\(F_{2}\) 分别是双曲线 \(C: \frac{x^{2}}{a^{2}} - \frac{y^{2}}{b^{2}} = 1 \) （\(a\)，\(b > 0\)）的左，右焦点，\(B\) 是虚轴的端点，直线 \(F_{1}B\) 与 \(C\) 的两条渐近线分别交于 \(P\)，\(Q\) 两点，线段 \(PQ\) 的垂直平分线与 \(x\) 轴交于点 \(M\). 若 \(|MF_{2}| = |F_{1}F_{2}|\)，则 \(C\) 的离心率是（\quad）

\choices
    {\(\frac{2\sqrt{3}}{3}\)}
    {\(\frac{\sqrt{6}}{2}\)}
    {\(\sqrt{2}\)}
    {\(\sqrt{3}\)}
\end{problem}

\begin{answer}
B
\end{answer}

\begin{solution}
令双曲线的半焦距为 \(c\)，则 \(c^2=a^2+b^2\)，且 \(F_1(-c,0)\)，\(F_2(c,0)\)，\(B(0,b)\). 直线 \(F_1B\) 上的点可表示为
\[
F_1+t(c,b)
\]
与渐近线 \(y=\frac ba x\) 相交时，代入参数式得
\[
tb=\frac ba(-c+tc)\Longrightarrow t(c-a)=c\Longrightarrow t_1=\frac{c}{c-a}
\]
与渐近线 \(y=-\frac ba x\) 相交时，代入参数式得
\[
tb=-\frac ba(-c+tc)\Longrightarrow t(c+a)=c\Longrightarrow t_2=\frac{c}{c+a}
\]
设 \(H\) 为线段 \(PQ\) 的中点，则
\[
H=F_1+\frac{t_1+t_2}{2}(c,b)
\]
又有
\[
\frac{t_1+t_2}{2}
=\frac c2\left(\frac1{c-a}+\frac1{c+a}\right)
=\frac{c^2}{c^2-a^2}
=\frac{c^2}{b^2}
\]
由于 \(PQ\) 的方向向量为 \((c,b)\)，其垂直平分线满足 \((X-H)\cdot(c,b)=0\). 设 \(M=(m,0)\)，则
\[
m=H_x+\frac bcH_y
=-c+\frac{c^2}{b^2}\left(c+\frac{b^2}{c}\right)
=\frac{c^3}{b^2}
\]
又因为 \(c^2>b^2\)，所以
\[
|MF_2|=m-c=\frac{ca^2}{b^2}
\]
由题意 \(|F_1F_2|=2c\)，故
\[
\frac{ca^2}{b^2}=2c\Longrightarrow a^2=2b^2
\]
因此离心率为
\[
e=\frac ca=\frac{\sqrt{a^2+b^2}}a=\sqrt{\frac32}=\frac{\sqrt6}{2}
\]
故选第二项.
\end{solution}

\begin{problem}
设 \(a > 0\)，\(b > 0\).（\quad）

\choices
    {若 \(2^{a} + 2a = 2^{b} + 3b\)，则 \(a > b\)}
    {若 \(2^{a} + 2a = 2^{b} + 3b\)，则 \(a < b\)}
    {若 \(2^{a} - 2a = 2^{b} - 3b\)，则 \(a > b\)}
    {若 \(2^{a} - 2a = 2^{b} - 3b\)，则 \(a < b\)}
\end{problem}

\begin{answer}
A
\end{answer}

\begin{solution}
令 \(g(x)=2^x+2x\). 因为
\[
g'(x)=2^x\ln2+2>0
\]
所以 \(g(x)\) 在 \((0,+\infty)\) 上严格递增. 若 \(2^a+2a=2^b+3b\)，则
\[
g(a)=2^b+3b>2^b+2b=g(b)
\]
从而 \(a>b\)，第一项正确，第二项错误.

对于后两个选项，先取 \(a=1\). 此时 \(2^a-2a=0\)，而函数 \(h(x)=2^x-3x\) 满足 \(h(3)=-1<0<h(4)=4\)，故由连续性存在 \(b\in(3,4)\) 使 \(h(b)=0\)，此时 \(a<b\)，第三项错误. 再令 \(a=b+1\)，方程 \(2^a-2a=2^b-3b\) 化为
\[
2^b=2-b
\]
函数 \(2^b+b-2\) 在 \(b=0\) 时为负，在 \(b=1\) 时为正，故存在 \(b\in(0,1)\) 满足该方程，此时 \(a>b\)，第四项也错误. 故选第一项.
\end{solution}

\begin{problem}
已知矩形 \(ABCD\)，\(AB = 1\)，\(BC = \sqrt{2}\)，将 \(\triangle ABD\) 沿矩形的对角线 \(BD\) 所在的直线进行翻折，在翻折过程中（\quad）

\choices
    {存在某个位置，使得直线 \(AC\) 与直线 \(BD\) 垂直}
    {存在某个位置，使得直线 \(AB\) 与直线 \(CD\) 垂直}
    {存在某个位置，使得直线 \(AD\) 与直线 \(BC\) 垂直}
    {对任意位置，三对直线“\(AC\) 与 \(BD\)”，“\(AB\) 与 \(CD\)”，“\(AD\) 与 \(BC\)”均不垂直}
\end{problem}

\begin{answer}
B
\end{answer}

\begin{solution}
取初始位置的坐标为
\(A(0,0,0)\)，\(B(1,0,0)\)，\(C(1,\sqrt2,0)\)，\(D(0,\sqrt2,0)\).
令两向量
\(\bs e=\frac1{\sqrt3}(-1,\sqrt2,0)\)，\(\bs n=\frac1{\sqrt3}(\sqrt2,1,0)\)
则 \(\bs e\) 是直线 \(BD\) 的单位方向向量，\(\bs n\) 与 \(\bs e\) 垂直. 设翻折后点 \(A\) 变为 \(A'\)，以折叠角为参数 \(\varphi\)，则
\[
\overrightarrow{BA'}=\frac1{\sqrt3}\bs e-\frac{\sqrt2}{\sqrt3}\bigl(\cos\varphi\,\bs n+\sin\varphi\,(0,0,1)\bigr)
\]
此外
\(\overrightarrow{BC}=\frac2{\sqrt3}\bs e+\frac{\sqrt2}{\sqrt3}\bs n\)，\(\overrightarrow{CD}=\frac1{\sqrt3}\bs e-\frac{\sqrt2}{\sqrt3}\bs n\)，\(\overrightarrow{BD}=\sqrt3\bs e\).
于是
\[
\overrightarrow{A'C}\cdot\overrightarrow{BD}
=(\overrightarrow{BC}-\overrightarrow{BA'})\cdot\sqrt3\bs e=1\neq0
\]
所以第一项所述位置不存在.

又有
\[
\overrightarrow{BA'}\cdot\overrightarrow{CD}=\frac13+\frac23\cos\varphi
\]
当 \(\cos\varphi=-\frac12\) 时该内积为零，故存在某个位置使 \(AB\perp CD\)，第二项正确.

最后
\[
\overrightarrow{DA'}\cdot\overrightarrow{BC}
=(\overrightarrow{BA'}-\overrightarrow{BD})\cdot\overrightarrow{BC}
=-\frac43-\frac23\cos\varphi
\]
要使其为零需 \(\cos\varphi=-2\)，不可能，因此第三项所述位置不存在，而第四项也不成立. 故选第二项.
\end{solution}

\section{填空题}

\begin{problem}
已知某三棱锥的三视图（单位：\(\mathrm{cm}\)）如图所示，则该三棱锥的体积等于\fillinblank{} \(\mathrm{cm}^3\).

\bitmapfigure[max width=0.78\linewidth]{img/2012/zhejiang_science/q11_fig1.png}
\end{problem}

\begin{answer}
1
\end{answer}

\begin{solution}
由三视图可取四个顶点的坐标为
\(O(0,0,0)\)，\(U(3,0,0)\)，\(V(3,1,0)\)，\(W(3,0,2)\).
其中底面 \(OUV\) 在水平面内，且
\[
S_{\triangle OUV}=\frac12\times3\times1=\frac32
\]
三棱锥的高为 \(W\) 到水平面的距离 \(2\). 因此体积为
\[
V=\frac13S_{\triangle OUV}\times2=1
\]
\end{solution}

\begin{problem}
若某程序框图如图所示，则该程序运行后输出的值是\fillinblank{}.

\begin{center}
\texfigure[max width=0.42\linewidth]{img_repaint/2012/zhejiang_science/q12_fig1.tex}
\end{center}
\end{problem}

\begin{answer}
\(\frac{1}{120}\)
\end{answer}

\begin{solution}
程序开始时 \(T=1\)，\(i=1\). 每次循环先执行 \(T=\frac Ti\)，再令 \(i=i+1\)，当 \(i>5\) 时输出. 因此依次除以 \(1\)、\(2\)、\(3\)、\(4\)、\(5\)，最后
\[
T=\frac1{1\times2\times3\times4\times5}=\frac1{120}
\]
\end{solution}

\begin{problem}
设公比为 \(q\)（\(q>0\)）的等比数列 \(\{a_{n}\}\) 的前 \(n\) 项和为 \(S_{n}\). 若 \(S_{2} = 3a_{2} + 2\)，\(S_{4} = 3a_{4} + 2\)，则 \(q =\fillinblank{}\).
\end{problem}

\begin{answer}
\(\frac{3}{2}\)
\end{answer}

\begin{solution}
设首项为 \(a_1\). 由 \(a_2=a_1q\) 及 \(S_2=a_1(1+q)\)，得
\[
a_1(1+q)=3a_1q+2\Longrightarrow a_1(1-2q)=2
\]
由第一式右端为 \(2\)，知 \(a_1\ne0\). 再由 \(a_4=a_1q^3\) 及 \(S_4=a_1(1+q+q^2+q^3)\)，得
\[
a_1(1+q+q^2+q^3)=3a_1q^3+2
\]
即得
\[
a_1(1+q+q^2-2q^3)=2
\]
两式相减并利用 \(a_1\ne0\)，得
\[
q+q^2-2q^3+2q=0
\Longrightarrow q(3+q-2q^2)=0
\]
由于 \(q>0\)，故 \(2q^2-q-3=0\)，从而 \(q=\frac32\).
\end{solution}

\begin{problem}
若将函数 \( f(x) = x^5 \) 表示为 \( f(x) = a_0 + a_1(1 + x) + a_2(1 + x)^2 + \cdots + a_5(1 + x)^5 \)，其中 \(a_0\)，\(a_1\)，\(a_2\)，\(\cdots\)，\(a_5\) 为实数，则 \( a_3 = \)\fillinblank{}.
\end{problem}

\begin{answer}
10
\end{answer}

\begin{solution}
令 \(u=1+x\)，则 \(x=u-1\). 由二项式定理，
\[
x^5=(u-1)^5=u^5-5u^4+10u^3-10u^2+5u-1
\]
与题设表示式逐项比较，\(u^3=(1+x)^3\) 的系数为 \(a_3=10\).
\end{solution}

\begin{problem}
在 \(\triangle ABC\) 中，\(M\) 是 \(BC\) 的中点，\(AM = 3\)，\(BC = 10\). 则 \(\overrightarrow{AB} \cdot \overrightarrow{AC} =\)\fillinblank{}.
\end{problem}

\begin{answer}
-16
\end{answer}

\begin{solution}
由于 \(M\) 是 \(BC\) 的中点，
\[
\overrightarrow{AM}=\frac{\overrightarrow{AB}+\overrightarrow{AC}}2
\]
所以
\[
|\overrightarrow{AB}+\overrightarrow{AC}|^2=4AM^2=36
\]
又有
\[
|\overrightarrow{AC}-\overrightarrow{AB}|^2=BC^2=100
\]
两式相减，得
\[
4\overrightarrow{AB}\cdot\overrightarrow{AC}=36-100=-64
\]
故 \(\overrightarrow{AB}\cdot\overrightarrow{AC}=-16\).
\end{solution}

\begin{problem}
定义：曲线 \(C\) 上的点到直线 \(l\) 的距离的最小值为曲线 \(C\) 到直线 \(l\) 的距离. 已知曲线 \( C_{1}: y = x^{2} + a \) 到直线 \(l: y = x\) 的距离等于曲线 \( C_{2}: x^{2} + (y + 4)^{2} = 2 \) 到直线 \(l: y = x\) 的距离，则实数 \( a = \)\fillinblank{}.
\end{problem}

\begin{answer}
\(\frac{9}{4}\)
\end{answer}

\begin{solution}
曲线 \(C_1\) 上点 \((x,x^2+a)\) 到直线 \(y=x\) 的距离为
\[
\frac{|x^2-x+a|}{\sqrt2}
\]
由于
\[
x^2-x=\left(x-\frac12\right)^2-\frac14
\]
且 \(x^2-x\) 的取值范围为 \(\left[-\frac14,+\infty\right)\). 若 \(a\leqslant\frac14\)，则存在 \(x\) 使 \(x^2-x+a=0\)，曲线 \(C_1\) 与直线 \(l\) 相交，二者距离为 \(0\)，不可能等于下述圆与直线的正距离. 因而 \(a>\frac14\)，从而 \(C_1\) 到直线 \(l\) 的距离为
\[
\frac{a-\frac14}{\sqrt2}
\]

曲线 \(C_2\) 是圆心 \((0,-4)\)、半径 \(\sqrt2\) 的圆. 圆心到直线 \(y=x\) 的距离为 \(2\sqrt2\)，故圆到该直线的最小距离为
\[
2\sqrt2-\sqrt2=\sqrt2
\]
因此
\[
\frac{a-\frac14}{\sqrt2}=\sqrt2
\Longrightarrow a=\frac94
\]
\end{solution}

\begin{problem}
设 \(a \in \R\)，若 \(x > 0\) 时均有 \([(a - 1)x - 1](x^2 - ax - 1) \geqslant 0\)，则 \(a = \)\fillinblank{}.
\end{problem}

\begin{answer}
\(\frac{3}{2}\)
\end{answer}

\begin{solution}
令函数
\(p(x)=(a-1)x-1\)，\(q(x)=x^2-ax-1\).
当 \(x\) 趋近于 \(0\) 且 \(x>0\) 时，\(p(x)<0\)，\(q(x)<0\)，故乘积为正. 多项式 \(q(x)\) 有唯一正根
\[
r=\frac{a+\sqrt{a^2+4}}2
\]
并且在 \((0,r)\) 上为负，在 \((r,+\infty)\) 上为正.

若 \(a\leqslant1\)，则对所有 \(x>0\) 有 \(p(x)<0\)，而当 \(x\) 充分大时 \(q(x)>0\)，这与乘积恒非负矛盾，故必须有 \(a>1\). 此时 \(p(x)\) 在正数范围内有唯一零点. 要使 \(p(x)q(x)\geqslant0\) 对所有正数 \(x\) 成立，\(p(x)\) 必须在 \((0,r)\) 上非正、在 \((r,+\infty)\) 上非负，因此 \(p(x)\) 的正根必须与 \(r\) 重合，即
\[
r=\frac1{a-1}
\]
将其代入 \(q(r)=0\)，得
\[
\frac1{(a-1)^2}-\frac a{a-1}-1=0
\Longrightarrow 1-a(a-1)-(a-1)^2=0
\Longrightarrow a(3-2a)=0
\]
结合 \(a>1\)，得到 \(a=\frac32\).

反之，当 \(a=\frac32\) 时，
\(p(x)=\frac{x-2}{2}\)，\(q(x)=x^2-\frac32x-1=(x-2)\left(x+\frac12\right)\)
所以
\[
p(x)q(x)=\frac12(x-2)^2\left(x+\frac12\right)\geqslant0
\]
对一切 \(x>0\) 成立，故所求为 \(a=\frac32\).
\end{solution}

\section{解答题}

\begin{problem}
在 \(\triangle ABC\) 中，内角 \(A\)，\(B\)，\(C\) 的对边分别为 \(a\)，\(b\)，\(c\). 已知 \(\cos A = \frac{2}{3}\)，\(\sin B = \sqrt{5} \cos C\).
\begin{enumerate}
\item 求 \(\tan C\) 的值；
\item 若 \(a = \sqrt{2}\)，求 \(\triangle ABC\) 的面积.
\end{enumerate}
\end{problem}

\begin{solution}
\begin{enumerate}
\item 由三角形内角和，\(B=\pi-A-C\)，所以
\[
\sin B=\sin(A+C)=\sin A\cos C+\cos A\sin C
\]
由 \(\cos A=\frac{2}{3}\) 及 \(0<A<\pi\)，得
\(\sin A=\frac{\sqrt{5}}{3}\). 代入已知条件，得
\[
\frac{\sqrt{5}}{3}\cos C+\frac{2}{3}\sin C=\sqrt{5}\cos C
\]
即 \(\sin C=\sqrt{5}\cos C\). 因为 \(0<C<\pi\)，且上式不可能有 \(\cos C=0\)，所以
\[
\tan C=\sqrt{5}
\]

\item 由 \(\tan C=\sqrt{5}\) 且 \(0<C<\pi\)，可知 \(C\) 为锐角，从而
\(\cos C=\frac{1}{\sqrt{6}}\)，\(\sin C=\frac{\sqrt{5}}{\sqrt{6}}\).
由正弦定理，
\[
c=a\frac{\sin C}{\sin A}=\sqrt{2}\cdot\frac{\sqrt{5}/\sqrt{6}}{\sqrt{5}/3}=\sqrt{3}
\]
又 \(\sin B=\sqrt{5}\cos C=\frac{\sqrt{5}}{\sqrt{6}}\)，故三角形面积为
\[
S_{\triangle ABC}=\frac{1}{2}ac\sin B
=\frac{1}{2}\cdot\sqrt{2}\cdot\sqrt{3}\cdot\frac{\sqrt{5}}{\sqrt{6}}
=\frac{\sqrt{5}}{2}
\]
\end{enumerate}
\end{solution}

\begin{problem}
已知箱中装有 \(4\) 个白球和 \(5\) 个黑球，且规定：取出一个白球得 \(2\) 分，取出一个黑球得 \(1\) 分. 现从该箱中任取（无放回，且每球取到的机会均等） \(3\) 个球，记随机变量 \(X\) 为取出此 \(3\) 个球所得分数之和.
\begin{enumerate}
\item 求 \(X\) 的分布列；
\item 求 \(E(X)\).
\end{enumerate}
\end{problem}

\begin{solution}
设取出的白球数为 \(k\)，则取出的黑球数为 \(3-k\)，从而
\(X=2k+(3-k)=3+k\)，\(k=0\)，\(1\)，\(2\)，\(3\).
从 \(9\) 个球中任取 \(3\) 个球的总取法数为
\(\mathrm C_9^3=84\). 因此
\[
P(X=3+k)=\frac{\mathrm C_4^k\mathrm C_5^{3-k}}{\mathrm C_9^3}
\]
具体计算得
\begin{center}
\begin{tabular}{c|cccc}
\(X\) & \(3\) & \(4\) & \(5\) & \(6\) \\
\hline
\(P\) & \(\frac{5}{42}\) & \(\frac{10}{21}\) & \(\frac{5}{14}\) & \(\frac{1}{21}\)
\end{tabular}
\end{center}

所以
\[
E(X)=3\cdot\frac{5}{42}+4\cdot\frac{10}{21}
+5\cdot\frac{5}{14}+6\cdot\frac{1}{21}
=\frac{13}{3}
\]
\end{solution}

\begin{problem}
如图，在四棱锥 \(P - ABCD\) 中，底面是边长为 \(2\sqrt{3}\) 的菱形，\(\angle BAD = 120^{\circ}\)，且 \(PA \perp\) 平面 \(ABCD\)，\(PA = 2\sqrt{6}\)，\(M\)，\(N\) 分别为 \(PB\)，\(PD\) 的中点.
\begin{enumerate}
\item 证明： \(MN \parallel\) 平面 \(ABCD\)；
\item 过点 \(A\) 作 \(AQ \perp PC\)，垂足为点 \(Q\)，求二面角 \(A - MN - Q\) 的平面角的余弦值.
\end{enumerate}
\bitmapfigure[max width=0.42\linewidth]{img/2012/zhejiang_science/q20_fig1.png}
\end{problem}

\begin{solution}
\begin{enumerate}
\item 在三角形 \(PBD\) 中，\(M\)，\(N\) 分别为 \(PB\)，\(PD\) 的中点，所以由中位线定理得
\(MN\parallel BD\). 又 \(BD\) 在平面 \(ABCD\) 内，而 \(P\notin\) 平面 \(ABCD\)，所以 \(M\notin\) 平面 \(ABCD\)，从而 \(MN\) 不在平面 \(ABCD\) 内，故
\[
MN\parallel\text{平面 }ABCD
\]

\item 建立如下坐标系：取
\(A(0,0,0)\)，\(B(2\sqrt{3},0,0)\)，\(D(-\sqrt{3},3,0)\)，\(C(\sqrt{3},3,0)\)，\(P(0,0,2\sqrt{6})\).
则 \(AB=AD=2\sqrt{3}\)，且 \(\angle BAD=120^\circ\)，符合题设. 由中点坐标公式，
\(M(\sqrt{3},0,\sqrt{6})\)，\(N\left(-\frac{\sqrt{3}}{2},\frac{3}{2},\sqrt{6}\right)\).
取 \(MN\) 的一个方向向量 \(\bs{u}=(-\sqrt{3},1,0)\). 又
\(\overrightarrow{PC}=(\sqrt{3},3,-2\sqrt{6})\)，\(|\overrightarrow{PC}|^2=36\).
因为 \(Q\) 在直线 \(PC\) 上且 \(AQ\perp PC\)，所以
\[
Q=P+\frac{(\overrightarrow{PA})\cdot\overrightarrow{PC}}{|\overrightarrow{PC}|^2}\overrightarrow{PC}
=\left(\frac{2\sqrt{3}}{3},2,\frac{2\sqrt{6}}{3}\right)
\]
平面 \(AMN\) 与平面 \(QMN\) 的法向量可分别取为
\[
\bs{n}_1=\bs{u}\times\overrightarrow{AM}
=(\sqrt{6},3\sqrt{2},-\sqrt{3})
\]
\[
\bs{n}_2=-3\bs{u}\times\overrightarrow{MQ}
=(\sqrt{6},3\sqrt{2},5\sqrt{3})
\]
于是二面角的平面角 \(\theta\) 的余弦为
\[
\cos\theta
=\frac{|\bs{n}_1\cdot\bs{n}_2|}{|\bs{n}_1||\bs{n}_2|}
=\frac{|6+18-15|}{\sqrt{27}\sqrt{99}}
=\frac{1}{\sqrt{33}}
=\frac{\sqrt{33}}{33}
\]
\end{enumerate}
\end{solution}

\begin{problem}
如图，椭圆 \(C: \frac{x^2}{a^2} + \frac{y^2}{b^2} = 1 \) （\(a > b > 0\)）的离心率为 \(\frac{1}{2}\). 其左焦点到点 \(P(2,1)\) 的距离为 \(\sqrt{10}\). 不经过原点 \(O\) 的直线 \(l\) 与 \(C\) 相交于 \(A\)，\(B\) 两点，且线段 \(AB\) 被直线 \(OP\) 平分.
\begin{enumerate}
\item 求椭圆 \(C\) 的方程；
\item 求 \( \triangle ABP \) 面积取最大值时直线 \(l\) 的方程.
\end{enumerate}
\bitmapfigure[max width=0.42\linewidth]{img/2012/zhejiang_science/q21_fig1.png}
\end{problem}

\begin{solution}
\begin{enumerate}
\item 设椭圆的半焦距为 \(c\). 由离心率
\[
\frac{c}{a}=\frac{1}{2}
\]
得 \(c=\frac{a}{2}\)，从而
\[
b^2=a^2-c^2=\frac{3a^2}{4}
\]
椭圆的左焦点为 \(F_1\left(-\frac{a}{2},0\right)\)，故由题意
\[
\left(2+\frac{a}{2}\right)^2+1=10
\]
因为 \(a>0\)，所以 \(2+\frac{a}{2}=3\)，即 \(a=2\). 因此
\(c=1\)，\(b^2=3\)，椭圆 \(C\) 的方程为
\[
\frac{x^2}{4}+\frac{y^2}{3}=1
\]

\item 设线段 \(AB\) 的中点为 \(M\). 由题意，\(M\) 在直线 \(OP\) 上；又因直线 \(l\) 不经过 \(O\)，故 \(M\neq O\). 于是可设 \(M=(2t,t)\)，其中 \(t\neq0\). 设直线 \(l\) 的一个方向向量为 \(\bs{d}=(u,v)\)，并写
\(A=M+\lambda\bs{d}\)，\(B=M-\lambda\bs{d}\)
其中 \(\lambda\neq0\). 由 \(A\)，\(B\) 同在椭圆上，代入椭圆方程并相减，得
\[
\frac{(2t+\lambda u)^2-(2t-\lambda u)^2}{4}+\frac{(t+\lambda v)^2-(t-\lambda v)^2}{3}=0
\]
即 \(2t\lambda u+\frac{4}{3}t\lambda v=0\). 由于 \(t\lambda\neq0\)，所以 \(3u+2v=0\). 因此直线 \(l\) 的方向可取 \((2,-3)\)，从而可设
\(3x+2y=k\)，\(k\neq0\).
直线 \(l\) 与 \(OP\) 的交点就是 \(M\)，故 \(M=\left(\frac{k}{4},\frac{k}{8}\right)\). 于是可设
\(A=M+\lambda(2,-3)\)，\(B=M-\lambda(2,-3)\).
代入椭圆方程，得
\[
\frac{(\frac{k}{4}+2\lambda)^2}{4}+\frac{(\frac{k}{8}-3\lambda)^2}{3}=1
\Longrightarrow 4\lambda^2+\frac{k^2}{48}=1
\]
所以 \(|k|<4\sqrt{3}\)，且
\[
|AB|=2|\lambda|\sqrt{13}=\sqrt{13}\sqrt{1-\frac{k^2}{48}}
\]
点 \(P(2,1)\) 到直线 \(l\) 的距离为
\[
d(P,l)=\frac{|8-k|}{\sqrt{13}}=\frac{8-k}{\sqrt{13}}
\]
其中最后一个等号由 \(k<4\sqrt{3}<8\) 得到. 因此
\[
S_{\triangle ABP}=\frac12|AB|d(P,l)=\frac12(8-k)\sqrt{1-\frac{k^2}{48}}
\]
令
\(g(k)=(8-k)^2\left(1-\frac{k^2}{48}\right)\)，\(-4\sqrt{3}<k<4\sqrt{3}\).
则
\[
g'(k)=\frac{(8-k)(k^2-4k-24)}{12}=\frac{(8-k)(k-2-2\sqrt{7})(k-2+2\sqrt{7})}{12}
\]
在区间 \((-4\sqrt{3},4\sqrt{3})\) 内，\(2+2\sqrt{7}>4\sqrt{3}\) 且 \(8-k>0\)，所以当 \(-4\sqrt{3}<k<2-2\sqrt{7}\) 时 \(g'(k)>0\)，当 \(2-2\sqrt{7}<k<4\sqrt{3}\) 时 \(g'(k)<0\). 因而 \(g(k)\) 在 \(k=2-2\sqrt{7}\) 处取得最大值. 由于 \(k\neq0\) 且 \(2-2\sqrt{7}\neq0\)，满足题设条件. 故所求直线方程为
\[
\boxed{3x+2y=2-2\sqrt{7}}
\]
\end{enumerate}
\end{solution}

\begin{problem}
已知 \(a > 0\)，\(b \in \R\)，函数 \(f(x) = 4ax^3 - 2bx - a + b\).
\begin{enumerate}
\item 证明：当 \(0 \leqslant x \leqslant 1\) 时，
\begin{enumerate}
\item 函数 \(f(x)\) 的最大值为 \(|2a-b|+a\)；
\item \(f(x)+|2a-b|+a\geqslant0\)；
\end{enumerate}
\item 若 \(-1 \leqslant f(x) \leqslant 1\) 对 \(x \in [0,1]\) 恒成立，求 \(a+b\) 的取值范围.
\end{enumerate}
\end{problem}

\begin{solution}
\begin{enumerate}
\item
\begin{enumerate}
\item 由
\[
f'(x)=12ax^2-2b=2(6ax^2-b)
\]
且 \(f''(x)=24ax\geqslant0\)（\(0\leqslant x\leqslant1\)），所以 \(f'(x)\) 在 \([0,1]\) 上单调不减. 若有内点零点，该点的导数符号只能由负变正，是极小值点；因此 \(f(x)\) 在 \([0,1]\) 上的最大值取在端点.
又有
\(f(0)=b-a\)，\(f(1)=3a-b\)
所以
\[
\max_{0\leqslant x\leqslant1}f(x)
=\max\{b-a,3a-b\}
=a+|2a-b|
\]

令 \(h(x)=2x^3-2x+1\). 由于
\(h'(x)=6x^2-2\)，\(h(x)\) 在 \([0,1]\) 上的最小值为
\[
h\left(\frac{1}{\sqrt{3}}\right)=1-\frac{4}{3\sqrt{3}}>0
\]
当 \(b\leqslant2a\) 时，
\[
f(x)+|2a-b|+a
=2(2ax^3-bx+a)
\geqslant2a(2x^3-2x+1)\geqslant0
\]
当 \(b\geqslant2a\) 时，
\[
f(x)+|2a-b|+a
=2\bigl(2ax^3+b(1-x)-a\bigr)
\geqslant2a(2x^3-2x+1)\geqslant0
\]
故结论成立.
\end{enumerate}

\item 若 \(-1\leqslant f(x)\leqslant1\) 在 \([0,1]\) 上恒成立，则由上一步的最大值公式，必有
\[
a+|2a-b|\leqslant1
\]
反之，若上式成立，则 \(f(x)\leqslant1\)，且由上一问
\[
f(x)\geqslant-\bigl(a+|2a-b|\bigr)\geqslant-1
\]
所以该条件也是充分的. 因而题设条件等价于
\[
a+|2a-b|\leqslant1
\]
由于 \(a>0\)，这等价于
\(0<a\leqslant1\)，\(3a-1\leqslant b\leqslant a+1\).
于是
\[
4a-1\leqslant a+b\leqslant2a+1
\]
从而 \(-1<a+b\leqslant3\). 下面说明区间内每个值都能取得. 设 \(s=a+b\)：
当 \(-1<s\leqslant1\) 时，取
\(a=\frac{s+1}{4}\)，\(b=s-a;\)
当 \(1<s\leqslant3\) 时，取
\(a=\frac{s-1}{2}\)，\(b=s-a\).
这两种取法均满足 \(0<a\leqslant1\) 及 \(3a-1\leqslant b\leqslant a+1\). 因此
\[
\boxed{-1<a+b\leqslant3}
\]
\end{enumerate}
\end{solution}
